\documentclass[lualatex, paper={90mm, 160mm}]{jlreq}
\usepackage{tikz}
\usepackage{lltjext}
\usetikzlibrary{calc}

\begin{document}
\pagestyle{empty} 

\begin{tikzpicture}[remember picture, overlay]
  % 1. 背景の原稿用紙風模様（アイデアA）
  \begin{scope}[opacity=0.2, color=black!60]
    % 縦線
    \foreach \x in {-4,...,4} {
      \draw[very thin] ($(current page.center) + (\x*8mm, -70mm)$) -- +(0, 140mm);
    }
    % 横線
    \foreach \y in {-8,...,8} {
      \draw[very thin] ($(current page.center) + (-32mm, \y*8mm)$) -- +(64mm, 0);
    }
    % 中央の飾り（原稿用紙の魚尾風）
    \node at ($(current page.center) + (0, 0)$) {\fontsize{8pt}{8pt}\selectfont $\langle \quad \rangle$};
  \end{scope}

  % 2. 表紙の二重枠線
  \draw[line width=1.2pt, color=black!80] ($(current page.south west) + (6mm, 6mm)$) rectangle ($(current page.north east) - (6mm, 6mm)$);
  \draw[line width=0.3pt, color=black!80] ($(current page.south west) + (7.5mm, 7.5mm)$) rectangle ($(current page.north east) - (7.5mm, 7.5mm)$);

  % 3. タイトル（縦書き）
  \node[anchor=north] at ($(current page.center) + (15mm, 50mm)$) {
    \begin{minipage}<t>{15em}
      \begin{center}
        {\fontsize{28pt}{40pt}\selectfont 走れメロス\par}
      \end{center}
    \end{minipage}
  };

  % 4. 作者名（縦書き）
  \node[anchor=north] at ($(current page.center) + (-15mm, -10mm)$) {
    \begin{minipage}<t>{10em}
      \begin{center}
        {\fontsize{16pt}{24pt}\selectfont 太宰 治\par}
      \end{center}
    \end{minipage}
  };
\end{tikzpicture}

\end{document}